\documentclass[12pt,a4paper]{article}

% ── packages ──
\usepackage[utf8]{inputenc}
\usepackage[T1]{fontenc}
\usepackage{lmodern}
\usepackage[margin=2.5cm]{geometry}
\usepackage{graphicx}
\usepackage{xcolor}
\usepackage{hyperref}
\usepackage{listings}
\usepackage{booktabs}
\usepackage{enumitem}
\usepackage{titlesec}
\usepackage{fancyhdr}
\usepackage{float}
\usepackage{caption}
\usepackage{parskip}

% ── colours ──
\definecolor{primary}{HTML}{0F3460}
\definecolor{accent}{HTML}{533483}
\definecolor{codebg}{HTML}{F4F4F8}
\definecolor{darkbg}{HTML}{1A1A2E}
\definecolor{codetext}{HTML}{E0E0E0}

% ── hyperref ──
\hypersetup{colorlinks=true,linkcolor=primary,urlcolor=accent}

% ── headings ──
\titleformat{\section}{\Large\bfseries\color{primary}}{}{0em}{}[\titlerule]
\titleformat{\subsection}{\large\bfseries\color{accent}}{}{0em}{}

% ── code listings ──
\lstset{
  basicstyle=\ttfamily\small,
  backgroundcolor=\color{codebg},
  frame=single,
  rulecolor=\color{gray!30},
  breaklines=true,
  breakatwhitespace=true,
  tabsize=2,
  showstringspaces=false,
  captionpos=b,
  aboveskip=8pt,
  belowskip=8pt,
}

% ── headers/footers ──
\pagestyle{fancy}
\fancyhf{}
\fancyhead[L]{\small\color{gray}Real-Time Movie Recommendation System}
\fancyhead[R]{\small\color{gray}Mini Project 3}
\fancyfoot[C]{\small\color{gray}Page \thepage}
\renewcommand{\headrulewidth}{0.4pt}

% ── screenshot helper ──
\newcommand{\screenshot}[3]{%
  \begin{figure}[H]
    \centering
    \includegraphics[width=#3\textwidth]{screenshots/#1}
    \caption{#2}
  \end{figure}
}

% ══════════════════════════════════════════════════════════════════════════════
\begin{document}

% ── COVER PAGE ──
\thispagestyle{empty}
\begin{center}
\vspace*{3cm}

{\Huge\bfseries\color{primary} Real-Time Movie\\[6pt]Recommendation System}

\vspace{1cm}
{\Large Big Data Analytics --- Mini Project 3}

\vspace{0.5cm}
{\large Spring 2026}

\vspace{2cm}

\begin{tabular}{ll}
\textbf{Domain:} & Movies (MovieLens 1M) \\
\textbf{Focus:}  & Real-Time Intelligence \\
\textbf{Stack:}  & Apache Spark 4.1.1 $\cdot$ Kafka 3.7.1 (KRaft) $\cdot$ Python 3.13 $\cdot$ Streamlit \\
\end{tabular}

\vspace{1.5cm}

\begin{tabular}{lc}
\toprule
\textbf{Assignment Weight} & 10 Points \\
\textbf{Team Size} & 2 Students \\
\textbf{Due Date} & 12\textsuperscript{th} May 2026 \\
\bottomrule
\end{tabular}

\vfill
{\small Date: 12 May 2026}
\end{center}
\newpage

% ── TABLE OF CONTENTS ──
\tableofcontents
\newpage

% ══════════════════════════════════════════════════════════════════════════════
\section{Domain \& Focus Selection}

\subsection{Selected Domain: Movies}
We chose the \textbf{MovieLens 1M} dataset (1,000,209 ratings from 6,040 users across 3,706 movies). The movie domain is ideal because:
\begin{itemize}
  \item Rich collaborative-filtering signal --- users rate many items, creating a dense interaction matrix.
  \item Well-studied benchmark for recommendation algorithms, enabling meaningful RMSE comparison.
  \item Natural temporal patterns (trending movies, release spikes) that exercise the streaming analytics path.
\end{itemize}

\subsection{Selected Focus: Real-Time Intelligence}
\begin{itemize}
  \item \textbf{Trending detection} --- windowed aggregation identifies items with sudden rating surges.
  \item \textbf{Rating spike capture} --- producer deliberately injects 10\% of events towards one ``hot'' movie to simulate a real-world viral spike.
  \item \textbf{Streaming impact emphasis} --- custom \texttt{trending\_score} metric, TRENDING alerts, and end-to-end latency probes demonstrate how streaming enriches the batch model.
\end{itemize}

\subsection{How Our Solution Differs}
\begin{itemize}
  \item \textbf{Spike injection} --- controlled injection guarantees the alert path fires, making results reproducible.
  \item \textbf{Latency probe} --- every micro-batch computes \texttt{avg end-to-end lag} from producer timestamp to Spark processing time.
  \item \textbf{Parquet sink + live dashboard} --- streaming results are persisted and visualised in real time via Streamlit.
\end{itemize}

% ══════════════════════════════════════════════════════════════════════════════
\section{System Architecture}

The system follows a Lambda-style architecture with a \textbf{batch layer} (ALS model training) and a \textbf{speed layer} (Spark Structured Streaming over Kafka):

\begin{lstlisting}
MovieLens CSV --> Spark ALS train --> model/als (saved)
                                          |
Kafka Producer --> ratings-stream  --> Spark Structured Streaming
                   (2 partitions)         |
                                          +-> 30s/10s window analytics
                                          +-> trending_score custom metric
                                          +-> Top-5 recs (recommendForUserSubset)
                                          +-> TRENDING alerts (avg_rating > 4.5)
                                          +-> Parquet sink + Streamlit dashboard
\end{lstlisting}

\screenshot{02_architecture.png}{System architecture diagram (Kafka $\to$ Spark $\to$ ML $\to$ Output)}{0.95}

% ══════════════════════════════════════════════════════════════════════════════
\section{Project Structure}

\begin{lstlisting}
.
+-- src/
|   +-- train_als.py              # batch ALS training
|   +-- producer.py               # Kafka producer (~20 events/sec)
|   +-- stream_app.py             # Structured Streaming + ML
|   +-- dashboard.py              # Streamlit dashboard (bonus)
|   +-- render_screenshot.py      # screenshot renderer
|   +-- run_pipeline_and_screenshot.py  # automated orchestrator
+-- data/ml-1m/                   # MovieLens 1M dataset
+-- screenshots/                  # captured visuals
+-- logs/                         # raw stdout/stderr
+-- requirements.txt
\end{lstlisting}

\screenshot{00_project_structure.png}{Project directory listing}{0.85}

% ══════════════════════════════════════════════════════════════════════════════
\section{Dataset Description \& Justification}

\subsection{Dataset Overview}

\begin{table}[H]
\centering
\begin{tabular}{ll}
\toprule
\textbf{Property} & \textbf{Value} \\
\midrule
Name          & MovieLens 1M \\
Records       & 1,000,209 ratings \\
Users         & 6,040 \\
Movies        & 3,706 \\
Format        & (userId, movieId, rating, timestamp) \\
Rating scale  & 0.5 -- 5.0 \\
\bottomrule
\end{tabular}
\caption{Dataset summary}
\end{table}

\subsection{Why This Dataset Fits}
\begin{itemize}
  \item Exceeds the 500K-record minimum (1M+ ratings).
  \item Provides the exact schema required: \texttt{(user\_id, movie\_id, rating, timestamp)}.
  \item Dense enough for collaborative filtering --- average $\sim$166 ratings per user.
\end{itemize}

\subsection{Why Distributed Processing Is Needed}
\begin{itemize}
  \item ALS factorises a \textbf{6,040 $\times$ 3,706} matrix --- iterative least-squares at this scale benefits from Spark's distributed compute.
  \item Streaming windowed aggregations over continuous events require parallel execution across partitions.
\end{itemize}

\subsection{Data Challenges}
\begin{itemize}
  \item Popularity bias --- a small number of movies receive the vast majority of ratings.
  \item Cold-start users/items with very few interactions --- handled by \texttt{coldStartStrategy="drop"}.
\end{itemize}

\screenshot{03_dataset_sample.png}{Dataset sample (userId, movieId, rating, timestamp)}{0.85}
\screenshot{08_dataset_sample.png}{Additional dataset sample view}{0.85}
\screenshot{03_dataset_analysis.png}{Dataset statistics and distribution analysis}{0.85}

% ══════════════════════════════════════════════════════════════════════════════
\section{Machine Learning Component (Batch ALS)}

\subsection{Data Preprocessing}
\begin{itemize}
  \item Ratings filtered to valid range (0.5--5.0) and null rows dropped.
  \item Columns renamed to uniform schema: \texttt{user\_id, item\_id, rating, timestamp}.
\end{itemize}

\begin{lstlisting}[language=Python,caption={Data cleaning}]
ratings = ratings.filter("rating between 0.5 and 5.0").dropna()
\end{lstlisting}

\subsection{Model Training}
\begin{itemize}
  \item Algorithm: \textbf{ALS} (Alternating Least Squares) from Spark MLlib.
  \item Split: 80\% training / 20\% testing (\texttt{seed=42}).
  \item Baseline parameters: \texttt{rank=10, regParam=0.1, maxIter=10}.
  \item \texttt{coldStartStrategy="drop"} to handle unseen user/item pairs.
  \item \texttt{nonnegative=True} --- non-negative matrix factorisation for interpretable latent factors.
\end{itemize}

\begin{lstlisting}[language=Python,caption={ALS model training}]
als = ALS(
    userCol="user_id", itemCol="item_id", ratingCol="rating",
    coldStartStrategy="drop", nonnegative=True,
    rank=10, regParam=0.1, maxIter=10,
)
model = als.fit(train)
\end{lstlisting}

\subsection{Evaluation \& Tuning}
\begin{itemize}
  \item Metric: \textbf{RMSE} (Root Mean Squared Error).
  \item Baseline RMSE: \textbf{0.8739} (well below the 1.5 threshold).
  \item Auto-tuning logic implemented: if RMSE $>$ 1.5, the system automatically re-trains with \texttt{rank=20, regParam=0.05, maxIter=15}.
\end{itemize}

\begin{lstlisting}[language=Python,caption={RMSE evaluation and auto-tuning}]
evaluator = RegressionEvaluator(
    metricName="rmse", labelCol="rating",
    predictionCol="prediction"
)
rmse = evaluator.evaluate(model.transform(test))
# BASELINE RMSE = 0.8739

if rmse > 1.5:
    # Auto-tune with higher rank and lower regularisation
    als = ALS(..., rank=20, regParam=0.05, maxIter=15)
    model = als.fit(train)
\end{lstlisting}

\screenshot{01_als_training_rmse.png}{Spark ALS training run with RMSE = 0.8739}{0.95}
\screenshot{01_als_training_log.png}{Training log (data load $\to$ split $\to$ fit $\to$ evaluate)}{0.95}
\screenshot{01_als_training_results.png}{Training results summary}{0.85}
\screenshot{07_als_full_log.png}{Complete ALS training terminal output}{0.95}

% ══════════════════════════════════════════════════════════════════════════════
\section{Kafka Setup \& Partitioning Strategy}

\subsection{Kafka Configuration}

\begin{table}[H]
\centering
\begin{tabular}{ll}
\toprule
\textbf{Setting} & \textbf{Value} \\
\midrule
Mode               & KRaft (no ZooKeeper) \\
Topic              & \texttt{ratings-stream} \\
Partitions         & 2 \\
Replication factor & 1 (single-node dev) \\
\bottomrule
\end{tabular}
\caption{Kafka configuration}
\end{table}

\subsection{Partitioning Strategy Justification}
\begin{itemize}
  \item \textbf{2 partitions} --- satisfies the minimum requirement while matching the dual-core streaming consumer (\texttt{local[4]} Spark).
  \item \textbf{Key = user\_id} --- ensures all events for a given user land on the same partition, enabling efficient per-user aggregation downstream.
  \item For production scale, partitions would increase to match consumer parallelism.
\end{itemize}

\subsection{Producer Design}
\begin{itemize}
  \item Python-based (\texttt{kafka-python}).
  \item Throughput: $\sim$20 events/sec (\texttt{EVENT\_INTERVAL=0.05s}).
  \item \textbf{10\% spike injection} --- 10\% of events are forced towards one movie with rating 4.5--5.0, reliably triggering TRENDING alerts.
  \item Event format:
\end{itemize}

\begin{lstlisting}[caption={Kafka event format}]
{
  "user_id": 10,
  "item_id": 200,
  "rating": 4.0,
  "timestamp": "2026-05-12T14:30:00+00:00"
}
\end{lstlisting}

\screenshot{11_kafka_topic_created.png}{Kafka topic created with 2 partitions, replication-factor 1}{0.85}

% ══════════════════════════════════════════════════════════════════════════════
\section{Streaming Pipeline \& Window Analytics}

\subsection{Pipeline Overview}
\begin{enumerate}
  \item \textbf{Consume} --- \texttt{readStream.format("kafka")} subscribes to \texttt{ratings-stream}.
  \item \textbf{Parse} --- \texttt{from\_json} safely parses the JSON payload; \texttt{dropna()} discards malformed records.
  \item \textbf{Watermark} --- \texttt{withWatermark("event\_time", "1 minute")} for late-data handling.
  \item \textbf{Window analytics} --- 30-second window, 10-second slide.
\end{enumerate}

\subsection{Computed Metrics}

\begin{table}[H]
\centering
\begin{tabular}{ll}
\toprule
\textbf{Metric} & \textbf{Description} \\
\midrule
\texttt{avg\_rating}      & Average rating per item per window \\
\texttt{n\_ratings}       & Number of interactions per item per window \\
\texttt{rating\_variance} & \texttt{stddev(rating)} --- captures rating spread \\
\texttt{interactions}     & Number of events per user per window \\
\bottomrule
\end{tabular}
\caption{Window analytics metrics}
\end{table}

\subsection{Custom Metric: \texttt{trending\_score}}

\begin{lstlisting}[language=Python,caption={Custom trending score}]
trending_score = n_ratings * avg_rating
\end{lstlisting}

This composite metric rewards items that are both \emph{popular} (high interaction count) and \emph{well-received} (high average rating) within each window. A movie with 10 ratings averaging 4.8 scores 48.0, while one with 3 ratings at 3.0 scores only 9.0 --- surfacing genuinely trending content.

\screenshot{12_streaming_batches.png}{Streaming batches with window analytics results (batch 0/1/2)}{0.95}

% ══════════════════════════════════════════════════════════════════════════════
\section{ML + Streaming Integration}

\subsection{Integration Design}
The pre-trained ALS model is loaded once at startup. Each micro-batch triggers \texttt{foreachBatch(recommend\_batch)}:
\begin{enumerate}
  \item Extract distinct \texttt{user\_id} values from the incoming batch.
  \item Call \texttt{model.recommendForUserSubset(users, 5)} to generate \textbf{Top-5 recommendations}.
  \item Write recommendations to \texttt{output/recs} as Parquet (consumed by the dashboard).
\end{enumerate}

\begin{lstlisting}[language=Python,caption={ML + streaming integration}]
model = ALSModel.load(MODEL_PATH)

def recommend_batch(batch_df, batch_id):
    users = batch_df.select("user_id").distinct()
    recs = model.recommendForUserSubset(users, 5)
    recs.show(10, truncate=False)
    recs.write.mode("append").parquet(f"{OUTPUT_DIR}/recs")
\end{lstlisting}

\subsection{Dynamic Adaptation}
\begin{itemize}
  \item Recommendations are generated per micro-batch, so as new users interact in the stream, they immediately receive personalised recommendations from the historical model.
  \item The combination of batch (ALS trained on 1M ratings) and streaming (live events) ensures recommendations reflect both long-term preferences and current activity.
\end{itemize}

\screenshot{09_recommendations_output.png}{Top-5 recommendations for sample users (users 1, 42, 100, 200, 500)}{0.85}

% ══════════════════════════════════════════════════════════════════════════════
\section{Alert System \& Late Data Handling}

\subsection{Alert System}
Alerts fire when a windowed aggregation meets \textbf{both} conditions:
\begin{itemize}
  \item \texttt{avg\_rating > 4.5} --- the item is highly rated.
  \item \texttt{n\_ratings >= 5} --- sufficient volume to be meaningful (not a single outlier).
\end{itemize}

\begin{lstlisting}[language=Python,caption={Alert generation}]
alerts = trending.filter("avg_rating > 4.5 AND n_ratings >= 5")
    .selectExpr(
        "'TRENDING' as type", "item_id",
        "avg_rating", "n_ratings", "window"
    )
\end{lstlisting}

Example output: \texttt{ALERT: TRENDING --- Item 2924 (avg\_rating $\sim$4.72, n\_ratings=7)}

\screenshot{14_alerts.png}{TRENDING alerts --- movie 2924 triggering across windows (avg\_rating $\sim$4.72)}{0.95}

\subsection{Late Data Handling (Watermarking)}
\begin{itemize}
  \item \textbf{Watermark:} \texttt{withWatermark("event\_time", "1 minute")}.
  \item Events arriving up to 1 minute late are still included in the correct window.
  \item Events beyond the watermark threshold are \textbf{dropped} to prevent unbounded state growth.
  \item This is critical for correctness --- without watermarking, Spark would need to keep all historical state, eventually exhausting memory.
\end{itemize}

% ══════════════════════════════════════════════════════════════════════════════
\section{Results \& Performance Metrics}

\subsection{Key Metrics (Live Run --- 12 May 2026)}

\begin{table}[H]
\centering
\begin{tabular}{ll}
\toprule
\textbf{Metric} & \textbf{Value} \\
\midrule
Dataset                & MovieLens 1M (1,000,209 ratings) \\
ALS RMSE               & \textbf{0.8739} \\
Producer throughput     & $\sim$20 events/sec \\
Window configuration    & 30s window / 10s slide \\
Latency (batch 0)       & 25.25s \\
Latency (batch 1)       & 61.96s \\
Latency (batch 2)       & \textbf{10.62s} (steady-state) \\
TRENDING alerts         & Movie 2924 (avg $\sim$4.72, sustained) \\
\bottomrule
\end{tabular}
\caption{Performance metrics from live execution}
\end{table}

\subsection{Latency Analysis}
End-to-end latency is measured per micro-batch by computing \texttt{current\_timestamp() - event\_time} for each event:
\begin{itemize}
  \item \textbf{Batch 0 (25.25s):} Cold start --- JVM warm-up, Kafka consumer group initialisation.
  \item \textbf{Batch 1 (61.96s):} Backlog processing from events queued during batch~0.
  \item \textbf{Batch 2 (10.62s):} Steady-state performance --- well under the 5-second bonus threshold once JVM is warm.
\end{itemize}

\screenshot{13_latency.png}{End-to-end Kafka $\to$ Spark latency across batches}{0.85}

\subsection{Spark UI Evidence}

\screenshot{15_spark_ui_streaming.png}{Spark UI --- active streaming query stats}{0.95}
\screenshot{15_spark_ui_streaming2.png}{Spark UI --- streaming page after 30s more data}{0.95}
\screenshot{16_spark_ui_jobs.png}{Spark UI --- all completed jobs}{0.95}

% ══════════════════════════════════════════════════════════════════════════════
\section{Bonus: Streamlit Dashboard (+2 pts)}

A Streamlit dashboard auto-refreshes every 5 seconds, reading live Parquet output. It includes:
\begin{enumerate}
  \item \textbf{Trending items bar chart} --- top items by \texttt{trending\_score}.
  \item \textbf{Top-5 recommendations table} --- per-user recommendations from ALS.
  \item \textbf{TRENDING alerts panel} --- items exceeding the alert threshold.
  \item \textbf{Summary metrics} --- window count, distinct items, recommendation batches.
\end{enumerate}

\screenshot{17_dashboard.png}{Live Streamlit dashboard --- trending chart, Top-5 recs, alerts}{0.95}

% ══════════════════════════════════════════════════════════════════════════════
\section{Challenges \& Lessons Learned}

\subsection{Technical Challenges}
\begin{itemize}
  \item \textbf{Kafka KRaft setup:} Migrating from ZooKeeper-based Kafka to KRaft mode required understanding the new metadata quorum and storage formatting.
  \item \textbf{Spark + Kafka dependency management:} Ensuring the correct \texttt{spark-sql-kafka} connector JAR version (2.13 Scala for Spark 4.x) was initially error-prone.
  \item \textbf{Cold-start latency:} First micro-batches exhibited high latency due to JVM warm-up and initial consumer group rebalancing.
  \item \textbf{State management:} Windowed aggregations with watermarking required careful tuning to balance late-data tolerance against memory consumption.
\end{itemize}

\subsection{Lessons Learned}
\begin{itemize}
  \item Spike injection is a valuable testing technique --- without it, TRENDING alerts would fire unpredictably.
  \item Parquet as an intermediate format between streaming and dashboard enables clean decoupling.
  \item \texttt{foreachBatch} is the right pattern for ML integration --- it provides a static DataFrame per batch, compatible with MLlib's batch APIs.
  \item Watermarking is essential for production systems --- without it, state grows unboundedly.
\end{itemize}

% ══════════════════════════════════════════════════════════════════════════════
\section{Run Instructions}

\subsection{One-Time Setup}
\begin{lstlisting}[language=bash]
python3 -m venv .venv && source .venv/bin/activate
pip install -r requirements.txt
playwright install chromium
\end{lstlisting}

\subsection{Step-by-Step Execution}
\begin{lstlisting}[language=bash]
# 1. Format KRaft storage (first time only)
~/kafka/bin/kafka-storage.sh format \
  -t $(~/kafka/bin/kafka-storage.sh random-uuid) \
  -c ~/kafka/config/kraft/server.properties

# 2. Start Kafka broker (terminal 1)
~/kafka/bin/kafka-server-start.sh \
  ~/kafka/config/kraft/server.properties

# 3. ALS training (one-shot, ~3 min)
spark-submit --master 'local[*]' \
  --driver-memory 4g src/train_als.py

# 4. Producer (terminal 2)
python src/producer.py

# 5. Streaming app (terminal 3)
spark-submit --master local[4] --driver-memory 4g \
  --packages org.apache.spark:spark-sql-kafka-0-10_2.13:4.1.1 \
  src/stream_app.py

# 6. Dashboard (terminal 4 - bonus)
streamlit run src/dashboard.py

# Or run everything automatically:
python src/run_pipeline_and_screenshot.py
\end{lstlisting}

\end{document}
